\scanpage{99}
\mspage{49}

Ceci posé, je dis que

\[
\begin{aligned}
 &\text{L'application (61) }\quad\operatorname{Et}_{P_{i_0}}(s)\overset{\lambda_s}{\overset{\sim}{\longrightarrow}}\operatorname{Et}_{P_{i_1}}(s)\\
 &\text{définie par }\lambda\text{ est un }\deleted{\text{\unsure{bis}}}\text{isomorphisme}\\
 &\text{de torseurs, pour un isomorphisme}\\
 &\text{des groupes d'opérateurs (uniquement}\\
 &\text{déterminé bien sûr)}\\
 &\deleted{\unsure{\gamma}}\lambda_{\Gamma,s}:\quad\Gamma_{i_0}\overset{\sim}{\underset{\unsure{1}}{\longrightarrow}}\Gamma_{i_1},
\end{aligned}\tag*{(63)}
\]

% slanted
\margin{à \uline{vérifier}}

Quand on a choisi des orientations de $P_{i_0}$ et de $P_{i_1}$, i.e. des origines de $\mu(i_0)$ et $\mu(i_1)$, alors $\lambda_{\Gamma,s}$ s'explicite par un isomorphisme de groupes

\[
\Lambda\overset{\sim}{\longrightarrow}\Lambda
\]

donc par un élément \inserted{$\gamma$} de $\Lambda^*$, \deleted{\unsure{point|pris}}\inserted{quand} on change \deleted{d'orientation}\inserted{d'origine} dans l'un des \deleted{\unsure{ensembles ordonnés}}\inserted{torseurs} $\mu(i_0)$, $\mu(i_1)$, \deleted{\unsure{le}} $\gamma$ est remplacé par $-\gamma$. Ainsi, on peut interpréter $\lambda_{\Gamma,s}$ comme un élément

\[
\lambda_{\Gamma,s}\in\bigl[\underbrace{\Lambda\wedge_{\mu_2}\bigl(\mu(i_0)\wedge\mu(i_1)\bigr)}_{\text{$\Lambda$-module inversible}\quad\Gamma_{i_0i_1}}\bigr]^{*}\tag*{(64)}
\]

ou, je le rappelle, $\unsure{i_1|\mu_1}=\lambda.i_0$.

On a d'autre part une \deleted{application}\inserted{\unsure{isom}} \deleted{canonique} canonique

\[
\begin{aligned}
 \Gamma_{i_0i_1}^{*}/\{\pm1\}&\simeq\Lambda^*/\pm1\\
 \unsure{\xi|\zeta}&\longmapsto \xi^{\unsure{\natural|\sharp|4}}
\end{aligned}\tag*{(65)}
\]
% archive sheet number: 98
